% Level 3 homework, Week 12 -- Showcase and Graduation
\homeworkweek{12}{Showcase and Graduation}{15--20 minutes. Your last homework of Level 3 -- congratulations!}
\namedateline

\begin{core}[A letter to future-you]
Write a short letter to yourself to open in one year. Include: what you built, the bug you are proudest of fixing, a piece of advice for when coding gets hard, and one thing you want to be able to build by then.
\answerlines{8}
\end{core}

\begin{challenge}[Keep coding: pick from the menu]
Choose one item and plan your first step (with a parent, offline-first):
\begin{itemize}[itemsep=1pt]
  \item \textbf{Version 2:} add a level editor, a second enemy type or a settings file to your project. Commit each step.
  \item \textbf{New algorithm:} write a maze generator using recursion, or sort a high-score table by two keys.
  \item \textbf{New data:} run a (consent-first, anonymous) family survey and chart it.
  \item \textbf{New AI idea:} train the monster sorter on your own made-up cards with three labels.
  \item \textbf{Teach it:} show a younger friend or relative how to draw a fractal tree.
\end{itemize}
My choice and first step: \answerlines{2}
\end{challenge}

\begin{creative}[Your developer badge]
Design a badge for yourself as a Level 3 graduate: a symbol that stands for your project, your nickname, and three skills you are proud of (e.g. \enquote{recursion}, \enquote{git}, \enquote{tests}).
\answerspace{4cm}
\end{creative}

\begin{selfcheck}
  \item I gave my 2-minute talk and explained one design decision.
  \item I gave kind, specific feedback to others.
  \item I know how to look for coding communities safely, together with a parent.
  \item I have a plan for what to learn next.
\end{selfcheck}
